\begin{table}[htb!]
\centering
\caption{Robustness of the FEMA damage-distribution comparison to
county size. Left: Spearman correlation of county WUA with FEMA-verified
housing damage (as in the main text). Right: the same correlation after
dividing both quantities by the county's residential exposure units, so
that damage intensities rather than totals are compared. Counties
without exposure-unit data are dropped from the normalised comparison.}
\label{tab:damage_normalized}
\begin{tabular}{lrrrrrr}
\hline
 & \multicolumn{3}{c}{County totals} & \multicolumn{3}{c}{Per exposure unit} \\
Storm & $n$ & $r_s$ & $p$ & $n$ & $r_s$ & $p$ \\
\hline
Katrina (2005) & 16 & $+0.34$ & 0.192 & 16 & $+0.44$ & 0.090 \\
Isaac (2012) & 1 & -- & -- & 1 & -- & -- \\
Harvey (2017) & 9 & $+0.80$ & 0.010 & 9 & $+0.72$ & 0.030 \\
Irma (2017) & 4 & $-0.60$ & 0.400 & 4 & $+1.00$ & -- \\
Michael (2018) & 29 & $+0.75$ & $<0.001$ & 29 & $+0.79$ & $<0.001$ \\
Laura (2020) & 12 & $+0.65$ & 0.022 & 12 & $+0.72$ & 0.008 \\
Ida (2021) & 18 & $+0.40$ & 0.097 & 18 & $+0.60$ & 0.008 \\
Ian (2022) & 8 & $+0.81$ & 0.015 & 8 & $+0.74$ & 0.037 \\
\hline
\end{tabular}
\end{table}
